%%%%%%%%%%%%%%%%%%%%%%%%%
\section{Introduction} 
%%%%%%%%%%%%%%%%%%%%%%%%%


%%%%%%%%%%%%%%%%%%%%%%%%%%%%%%%%%%
\subsection{Purpose}
%%%%%%%%%%%%%%%%%%%%%%%%%%%%%%%%%%


This \ac{SDD} establishes the architectural framework for the ALMA \ac{WSU} at the preliminary design level. It describes the major components of the system (subsystems) and references detailed subsystem designs where available.

This document serves as a system-level technical reference, providing engineers, developers, project managers, and scientists with a comprehensive overview of the upgraded observatory architecture. The current version supports preparation for the  System \ac{PDR} scheduled for July 2025 and will evolve iteratively as the design matures through subsequent development phases.
This SDD serves as an explanatory snapshot of the design, with the binding requirements codified and maintained in the System Technical Requirements.

 
%%%%%%%%%%%%%%%%%%%%%%%%%%%%%%%%%%
\subsection{Scope}
%%%%%%%%%%%%%%%%%%%%%%%%%%%%%%%%%%

The document defines the  WSU system as envisioned after completion of all the WSU Program Milestones defined in the ALMA WSU Implementation Plan \cite{imp_plan_b}.
%The full system includes:
%\begin{itemize}[itemsep=-0.5em]
%    \item Deployment of wideband receivers for Bands 2, 6, 7, and 8.
%    \item A fourfold increase in bandwidth (32~GHz per polarization) with the \ac{ATAC} and the \ac{TPGS}.
%    \item Completion of the computing and data processing infrastructure.
%\end{itemize}

The scope of this document does not include
\begin{itemize}[itemsep=-0.5em]
    \item Planning and execution of \ac{AIVC} activities described in \cite{aivc_report}.
    \item Transition activities associated with Milestones $1-5$, defined in Section~\ref{intro:background} below and in \cite{imp_plan_b}, including the transition to WSU operations \cite{TWO} and the data processing transition \cite{dp_tp}.
    \item Systems and upgrades scheduled for implementation beyond Milestone 5.
    \item Developments in the ALMA Development Roadmap~\cite{road_map} outside the scope of the WSU.
\end{itemize}

The \ac{SDD} provides a high-level description of the WSU system. It does not present detailed subsystem designs, which will be addressed in separate design documents. The design presented herein is presented to satisfy the requirements~\cite{pa_req} for the System \ac{PDR} in July 2025. The \ac{SDD} will continue to evolve in accordance with the product design process defined in~\cite{pa_req}. 



%%%%%%%%%%%%%%%%%%%%%%%%%%%%%%%%%%
 \subsection{Background}
%%%%%%%%%%%%%%%%%%%%%%%%%%%%%%%%%%
\label{intro:background}

Since beginning science operations in 2011, ALMA has transformed our understanding of the cold universe. With unparalleled sensitivity and angular resolution at millimeter and submillimeter wavelengths, ALMA has enabled landmark discoveries in planetary system formation, the structure and chemistry of the interstellar medium, the birth of stars, and the evolution of galaxies across cosmic time. Reflecting its scientific impact, ALMA receives about 1700 observing proposals each year, with an oversubscription rate of approximately 7:1, making it one of the most competitive observatories in the world.

To meet growing scientific demand and unlock new discovery space, ALMA must overcome limitations in instantaneous bandwidth and spectral agility that constrain both the efficiency and breadth of current observations. The \ac{WSU}, identified as the top priority in the ALMA Development Roadmap \cite{road_map}, addresses these challenges by significantly expanding spectral coverage and improving sensitivity \cite{memo621}. This comprehensive system-wide upgrade will enable a fourfold increase in instantaneous bandwidth for both interferometric and single-dish (Total Power) observations. Achieving this capability requires substantial upgrades across multiple subsystems, including:
\begin{itemize}[itemsep=-0.5em]
    \item Wideband receivers with 16~GHz per polarization per sideband (32~GHz total observing bandwidth per polarization), which is $4\times$ the current bandwidth of most ALMA receivers.
    \item New digitizers, associated \ac{DTS} components, and fiber optics system.
    \item A new correlator for interferometric and \ac{VLBI}/Phased Array (a.k.a., Beamforming) observations.
    \item A new spectrometer for single-dish (Total Power) observations.
    \item Enhancements to the computing infrastructure and data processing infrastructure.
\end{itemize}
These upgrades will provide ALMA with the capacity to handle larger data volumes and offer greater flexibility for observing a wider range of scientific targets. % , ensuring that ALMA remains at the forefront of astronomical research.

The WSU program is being executed under a structured systems engineering framework that ensures traceability from scientific goals to system-level requirements and technical implementation. The process began with the development of top-level Science Requirements~\cite{sci_specs}, which served as the basis for the \ac{SRR}. The \ac{SRR} confirmed the scientific priorities and established a firm foundation for design.

Building on that foundation, the \ac{CoSDD} refined the system architecture and defined critical subsystem interfaces through iterative development. This \ac{SDD} advances the architectural baseline to the preliminary design level and supports the System \ac{PDR} scheduled for July 2025. 
%The \ac{PDR} will assess whether the design is sufficiently mature to meet the performance, schedule, and integration goals required to realize the WSU.

% System and subsystem-level oversight of the WSU is managed by the \ac{AMT}. The \ac{ADD-Dev} is the WSU Project Director and is responsible for the planning and implementation of the WSU, which includes coordinating activities across ALMA in order to ensure timely delivery, integration, verification, and compliance of the entire WSU system. The \ac{JDT} supports the \ac{ADD-Dev} in this role and serves as the WSU Project Office, building and monitoring system-level schedules, ensuring project management and system engineering processes are in place for the WSU, and managing onsite implementation.

System and subsystem-level oversight of the WSU is managed by the \ac{AMT}. The \ac{ADD-Dev}, as WSU Project Director, is responsible for system-level planning and implementation, including coordinating activities across ALMA to ensure timely delivery, integration, verification, and compliance of the full WSU system. The \ac{JDT} supports the \ac{ADD-Dev} in this role and serves as the WSU Project Office, building and monitoring system-level schedules, ensuring project management and system engineering processes are in place, and managing onsite implementation. The Executive's members of the AMT, the Executive Program(me) Managers, are responsible for  subsystems delivered by their region's Development Projects.  Each subsystem undergoes its own design, development, and review process. Table~\ref{tbl:status_projects} summarizes the status of individual subsystems.

\begin{table}[htb!]
\centering
\caption{Status of WSU subsystem development projects.}
\label{tbl:status_projects}
\begin{tabularx}{\textwidth}{|l|X|X|}
\hline
\textbf{Project} & \textbf{Status} & \textbf{Next milestone}\\
\hline
Band 2 & Deployment underway & Available starting Cycle 13 (October 2026)\\
\hline
Band 6v2 & Prototype under development & PDR in Q4 2025\\
\hline
Band 7v2 & In study phase & \\
\hline
Band 8v2 & Prototype under development & PDR in Q1-Q2 2026\\
\hline
WIFP & PDR conditionally passed (November 2024) & \acs{CDR} in Q4 2026\\
\hline
DTS & \acs{PDR} passed (October 2024) & \acs{CDMR} in Q3-Q4 2026\\
\hline
FOS & Project approved by ALMA Board (April 2025)  & \acs{CDMR} in Q4 2025\\
\hline 
ATAC & \ac{PDR} passed (June 2024) & \acs{CDR} in Q4 2025\\
\hline
TPGS & \acs{CoDR}/SRR passed (November 2024) & \acs{PDR} in Q4 2025\\
\hline
OCRO & In detailed design phase & \acs{CDMR} in Q2 2025\\
\hline
RADPS-ALMA & RADPS SRR passed (February 2025) & RADPS CoDR in Q3 2025, RADPS-ALMA \acs{PDR} (mid-2026)\\
\hline
WSU Software & In Preliminary Design phase & PDR in March 2026\\
\hline
\end{tabularx}
\end{table}

Deployment of the WSU will follow a phased approach as defined in the WSU Implementation Plan~\cite{imp_plan_b}. The process is structured around five milestones that mark key transition points, from initial deployment to full scientific capability, and guide the integration, testing, and commissioning of upgraded components across the array. The major components of each milestone are as follows:

% Milestone 1
\noindent\textcolor{milestoneblue}{\textbf{Milestone 1: Initial WSU Scientific Observations}}\par
\vspace{-0.6\baselineskip} % Pulls bullets up close to title
\begin{itemize}[itemsep=-0.5em]
    \item Sensitivity similar to the legacy system; top-priority modes commissioned.
    \item Legacy receivers for Bands 1/3--10; WSU Band 2 receiver with 4$\times$ IF bandwidth.
    \item At least $36\times12$-m antennas retrofitted with \ac{WIFP} and \ac{DTS}. 
    \item \ac{ATAC} supports 4$\times$ ingest, 2$\times$ correlation.
    \item Initial software and archive capabilities with  limits on data rates and size of data products.
    \item Legacy pipeline adapted for WSU calibration and imaging integrated into RADPS-ALMA  workflow and resource manager (State 4 in data processing transition plan~\cite{dp_tp}).
\end{itemize}\vspace{0pt}

% Milestone 2
\noindent\textcolor{milestoneblue}{\textbf{Milestone 2: End of System \ac{AIV}}}\par
\vspace{-0.6\baselineskip}
\begin{itemize}[itemsep=-0.5em]
    \item All antennas retrofitted with \ac{WIFP} and \ac{DTS}.
    \item \ac{TPGS} \ac{AIV} completed.
\end{itemize}\vspace{0pt}

% Milestone 3
\noindent\textcolor{milestoneblue}{\textbf{Milestone 3: End of Commissioning}}\par
\vspace{-0.6\baselineskip}
\begin{itemize}[itemsep=-0.5em]
    \item All legacy ALMA observing modes commissioned with WSU hardware at 2$\times$ bandwidth.
\end{itemize}\vspace{0pt}

% Milestone 4
\noindent\textcolor{milestoneblue}{\textbf{Milestone 4: End of Data Processing Transition}}\par
\vspace{-0.6\baselineskip}
\begin{itemize}[itemsep=-0.5em]
    \item Full \ac{WSU} data processing software deployed capable of handling $4\times$ bandwidth.
    \item Restrictions lifted on the allowed data rate and the size of science products.
\end{itemize}\vspace{0pt}

% Milestone 5
\noindent\textcolor{milestoneblue}{\textbf{Milestone 5: Top Scientific Priorities Achieved}}\par
\vspace{-0.6\baselineskip}
\begin{itemize}[itemsep=-0.5em]
    \item New Bands 6, 7, and 8 wideband receivers installed.
    \item \ac{ATAC} expanded to process 32~GHz/polarization (4$\times$ the current bandwidth).
    \item Final computing and storage capacity in place.
\end{itemize}

The design presented in this \ac{SDD} corresponds to the system configuration at the completion of Milestones $1-5$, which marks the completion of major upgrades and the realization of ALMA’s top scientific priorities for the WSU.

%%%%%%%%%%%%%%%%%%%%%%%%%%%%%%%%%%
\subsection{Document Organization}
%%%%%%%%%%%%%%%%%%%%%%%%%%%%%%%%%%

This document presents the WSU system design, structured into five main parts.

\begin{enumerate}
\item \textbf{Processes}\\
Section~\ref{sec:wsu_processes} outlines the end-to-end operational workflows, from proposal submission to data delivery. These workflows are based on the current ALMA model, refined over more than a decade of operations. The section highlights where existing processes are retained and which modifications are required to support the upgraded system.

\item \textbf{Signal Chain}\\
Section~\ref{sec:wsu_signal_chain} details the hardware infrastructure of the \ac{WSU} signal path, including upgrades to all major components—from front-end receivers and digitizers to the data transmission system, correlator, and spectrometer. The section emphasizes the integrated design required to achieve the fourfold increase in bandwidth and maintain system-wide compatibility.

\item \textbf{Computing Infrastructure – IT}\\
Section~\ref{sec:wsu_compu_it} describes the computing architecture that supports data processing, storage, and archival services. It includes the transition to a new processing framework and the expansion of real-time and offline capabilities to accommodate the increased data volume.

\item \textbf{Computing Infrastructure - Software}\\
Section~\ref{sec:wsu_compu_sw} describes how the software will need to evolve to handle the increased data volume obtained with the WSU, including data processing.

\item \textbf{Design Functional and Non-Functional Constraints}\\
Section~\ref{sec:nf_design_constraints} describes the main functional and non-functional constraints (i.e., reliability, availability, maintainability, safety, power standards, etc.) that the WSU design is subject to, both for signal chain, computing equipment, and computing software.

\end{enumerate}

This document establishes the first baseline for the WSU System Design, covering the upgraded observatory processes, signal chain, computing infrastructure, and software components. The design builds on the existing ALMA system architecture, incorporating lessons learned from more than a decade of operations. Readers seeking additional context on the current ALMA architecture may refer to the companion document, \textit{Wideband Sensitivity Upgrade – Current ALMA System Architecture and the Associated WSU Impact Assessment (CSA-IA)} \cite{CSA-IA}, 
which includes an overview of the ALMA Telescope (CSA-IA Section~2),
a description of the current ALMA processes, signal chain, \ac{IT} infrastructure, and software infrastructure  (CSA-IA Sections $3-6$), and a WSU impact assessment for each of these areas (CSA-IA Sections $7-10$).




%%%%%%%%%%%%%%%%%%%%%%%%%%%%%%%%%%
\subsection{Applicable Documents}
%%%%%%%%%%%%%%%%%%%%%%%%%%%%%%%%%%
The following documents are part of this document to the extent specified herein. If not explicitly stated otherwise, the latest release of the document is applicable.


\vspace*{0.8cm}
%%%%%%%%%%%%%%%%%
\begin{center}
\newrefcontext[labelprefix=AD]
\printbibtabular[env=bibtabular, heading=none, keyword=AD] % Filter for AD
\end{center}
%%%%%%%%%%%%%%%%%

%%%%%%%%%%%%%%%%%%%%%%%%%%%%%%%%%%
\subsection{Reference Documents}
%%%%%%%%%%%%%%%%%%%%%%%%%%%%%%%%%%
The following documents contain additional information.

\vspace*{0.8cm}
%%%%%%%%%%%%%%%%%
\begin{center}
\newrefcontext[labelprefix=RD]
\printbibtabular[env=bibtabular, heading=none, keyword=RD] % Filter for RD
%\printbibtabular[env=bibtabular, heading=none, keyword=RD] % Filter for RD
\end{center}
%%%%%%%%%%%%%%%%%

% This is needed since the AD/RD lists are formally tables. This subtracts 2 fromthe table counter to ignore these tables. The other way to do this is use longtable* when generating the AD/RD tables, but this broke the link from the references cited in the text to the AD/RD tables.
\addtocounter{table}{-2}

\newpage

%%%%%%%%%%%%%%%%%%%%%%%%%%%%%%%%%%
\subsection{List of Acronyms}
\begin{acronym}[RADPS-ALMA]\itemsep0pt
%%%%%%%%%%%%%%%%%%
\acro{2SB}{Dual Sideband}
\acro{ABM}{Antenna Bus Master}
\acro{ACA}{Atacama Compact Array}
\acro{ACAC}{Atacama Compact Array Correlator}
\acro{ACASPEC}{Atacama Compact Array Spectrometer}
\acro{ACD}{Amplitude Calibration Device}
\acro{ACRV}{Acceptance Review}
\acro{ACS}{ALMA Common Software}
\acro{ACSE}{ALMA Computing Simulation Environment}
\acro{ACSOS}{Array Calibration and Observing Strategies}
\acro{ACU}{Antenna Control Unit}
\acro{AD}{Applicable Document}
\acro{ADAPT}{ALMA Data Processing Toolkit}
\acro{ADC}{ALMA Department of Computing}
\acro{ADD-Dev}{ALMA Deputy Director for Development}
\acro{ADD-Ops}{ALMA Deputy Director for Operations}
\acro{ADE}{ALMA Department of Engineering}
\acro{ADMIT}{ALMA Data Mining Toolkit}
\acro{AE}{Array Element}
\acro{AEDM}{ALMA Electronic Document Management System}
\acro{AEI}{Array Elements Integration}
\acro{AEIT}{Array Elements Integration Team}
\acro{AFSP}{ALMA Frequency Slice Processor}
\acro{AIV}{Assembly, Integration and Verification}
\acro{AIVC}{Assembly, Integration, Verification and Commissioning}
\acro{ALMA}{Atacama Large Millimetre/submillimetre Array}
\acro{AMAPOLA}{Analytic Matrix for ALMA POLArimetry}
\acro{AMC}{Active Multiplier Chain}
\acro{AMG}{Array Maintenance Group}
\acro{AMT}{ALMA Management Team}
\acro{AoD}{Astronomer on Duty}
\acroplural{AoD}[AoDs]{Astronomers on Duty}
\acro{AOG}{Array Operations Group}
\acro{AOPG}{Array Operations Performance Group}
\acro{AOS}{Array Operations Site}
\acro{APA}{AQUA Pipeline Agent}
\acro{APC}{Atmospheric Phase Correction}
\acro{APDM}{ALMA Project Data Model}
\acro{APE}{ALMA Production Environment}
\acro{APE-HILSE}{ALMA Production Environment for Hardware In the Loop Simulation Environment}
\acro{APO}{Archive and Pipeline Operations}
\acro{APRC}{ALMA Proposal Review Committee}
\acro{APS}{ALMA Phasing System}
\acro{AQ}{Archive Query}
\acro{AQUA}{ALMA Quality Assurance}
\acro{ARC}{ALMA Regional Center}
\acro{ARM}{Array Runtime Manager}
\acro{ARTM}{Array Real Time Machine}
\acro{ASA}{ALMA Science Archive}
\acro{ASC}{ALMA Support/Science Center}
\acro{ASDM}{ALMA Science Data Model}
%\acro{ASIAA}{}
\acro{ASIC}{Application Specific Integrated Circuit}
\acro{ASTR}{ALMA System Technical Requirements}
\acro{AT}{Array Time}
\acro{ATAC}{Advanced Technology ALMA Correlator}
\acro{ATAM}{Architecture Tradeoff Analysis Model}
\acro{AUI}{Associated Universities Incorporated}
\acro{AVCC}{ALMA Very Coarse Channelizer}
\acro{AZ}{Azimuth}
\acro{B2B}{Band-to-Band}
\acro{BACI}{Basic Access Control Interface}
\acro{BDF}{Binary Data Format}
\acro{BE}{Back-End}
\acro{BER}{Bit Error Rate}
\acro{BiCMOS}{Bipolar CMOS}
\acro{BLC}{ALMA Baseline Correlator}
\acro{BMS}{Building Management System}
\acro{BOM}{Bill Of Materials}
\acro{BTU}{British Thermal Unit}
\acro{BWSW}{Bandwidth Switching}
\acro{CAN}{Control Area Network}
\acro{CAPEX}{Capital Expenditures}
\acro{CARTA}{Cube Analysis and Rendering Tool for Astronomy}
\acro{CASA}{Common Astronomy Software Applications}
\acro{CBF}{Correlator Beam Former}
\acro{CBW}{Correlated Bandwidth}
\acro{CCA}{Cold Cartridge Assembly}
\acro{CCB}{Change Control Board}
\acro{CCL}{Control Command Language}
\acro{CCTL}{Commissioning Campaign Test Lead}
\acro{CDB}{Configuration Database (for ACS components)}
\acro{CDL}{Central Development Lab}
\acro{CDMR}{Critical Design and Manufacturing Review}
\acro{CDP}{Correlator Data Processor}
\acro{CDR}{Critical Design Review}
\acro{CDU}{Coolant Distribution Unit}
\acro{CFP}{C form-factor pluggable}
\acro{CfP}{Call for Proposals}
\acro{CLAI}{Cluster/Adapt Interface}
\acro{CLNA}{Cryogenic Low-Noise Amplifiers}
\acro{CLO}{Central Local Oscillator}
\acro{CLOA}{Central Local Oscillator Article}
\acro{CM}{Compact 7-m Array Mitsubishi antenna}
\acro{CMMS}{Computerized Maintenance Management System}
\acro{CMOS}{Complementary Metal-Oxide-Semiconductor}
\acro{CoDR}{Conceptual Design Review}
\acro{CORBA}{Common Object Request Broker Architecture}
\acro{CorrGUI}{Correlator Graphical User Interface}
\acro{CoSDD}{Conceptual System Design Description}
\acro{COTS}{Commercial Off-The-Shelf}
\acro{CPU}{Central Processing Unit}
\acro{CRAH}{Computer Room Air Handler}
\acro{CRD}{Central Reference Distributor}
\acro{CRE}{Change Request}
\acro{CRR}{Control Room Relocation}
\acro{CSA}{Current System Architecture}
\acro{CSA-IA}{Current System Architecture-Impact Assessment}
\acro{CSCG}{Control System Coordination Group}
\acro{CSV}{Commissioning and Science Verification}
\acro{CTG}{Commissioning Testing Group}
\acro{CU}{Control Unit}
\acro{CVR}{Central Variable Reference}
\acro{CVS}{Concurrent Versions System}
\acro{DA}{12-m Array Alcatel-ElE-MAN antenna}
\acro{DC}{Deployment Concept}
\acro{DCO}{Digital Coherent Optics}
\acro{DDR}{Double Data Rate Synchronous Dynamic Random Access Memory}
\acro{DDS}{Data Distribution Service}
\acro{DDT}{Director's Discretionary Time}
\acro{DFF}{D-type Flip-Flop}
\acro{DGC}{Differential Gain Calibrator}
\acro{DGCK}{Digitizer Clock}
\acro{DGCTR}{Digitizer Clock and Time Reference}
\acro{DGM}{Digitizer Module}
\acro{DGOT} {Digitizer Optical Transmission}
\acro{DMG}{Data Management Group}
\acro{DMS}{Data Management System (NRAO Department)}
\acro{DP}{Data Processing}
\acro{DPDA}{Data Processing, Distribution, and Accessibility}
\acro{DPR}{Distributed Peer Review}
%\acro{DPS}{Data Processing System}
\acro{DR}{Data Reducer}
\acro{DRRUP}{Data Rate Ramp-Up Plan}
\acro{DSA}{Dynamic Scheduler Algorithm}
\acro{DSB}{Double Sideband}
\acro{DSO}{Department of Science Operations}
\acro{DSPOT}{Digital Signal Processing and Optical Transmission}
\acro{DTAS}{Data Transfer and Archive Storage}
\acro{DTS}{Data Transmission System}
\acro{DTX}{Data Transmitter Line Replacement Unit}
\acro{DV}{12-m Array Vertex antenna}
\acro{DVM}{Digital Voltage Meter}
\acro{DWDM}{Dense Wavelength Division Multiplexing}
\acro{E2C}{Ethernet to CAN}
\acro{EA}{East Asian}
\acro{EB}{Execution Block}
\acro{EDFA}{Erbium-Doped Fiber Amplifier}
\acro{EDM}{Electronic Documentation Management}
\acro{EF}{Execution Fraction}
\acro{EHT}{Event Horizon Telescope}
\acro{EIA}{Electronic Industries Alliance}
\acro{ELT}{Extremely Large Telescope}
\acro{ELV}{Elevation}
\acro{ENOB}{Effective Number Of Bits}
\acro{EOC}{Extension and Optimization of Capabilities}
\acro{EOL}{End Of Life}
\acro{ESG}{Engineering Services Group}
\acro{ESN}{Electronic Serial Number}
\acro{ESO}{European Southern Observatory}
\acro{EU}{European}
\acro{F2F}{Face-to-face}
\acro{FDM}{Frequency Division Mode}
\acro{FE}{Front-End}
\acro{FEC}{Forward Error Correction}
\acro{FEMC}{Front-End Monitor and Control module}
\acro{FEPS}{Front-End Power Supply}
\acro{FESV}{Front-End Service Vehicle}
\acro{FETMS}{Front-End Test and Measurement System}
\acro{FFT}{Fast Fourier Transform}
\acro{FFX}{Frequency division / Fourier transform (or fine channelization) / Cross correlation}
\acro{FITS}{Flexible Image Transport System}
\acro{FLOOG}{First LO Offset Generator}
\acro{FLOP}{Floating-Point Operations}
\acro{FMEA}{Failure Mode and Effects Analysis}
\acro{FMECA}{Failure Mode, Effects and Criticality Analysis}
\acro{FO}{Fiber Optics}
\acro{FOS}{Fiber Optic System}
\acro{FoV}{Field of View}
\acro{FPGA}{Field Programmable Gate Array}
\acro{FS}{Frequency Slice}
\acro{FTE}{Full Time Equivalent}
\acro{FW}{Firmware}
\acro{FXF}{Frequency division / Cross correlation / Fourier transform}
\acro{Gb}{Gigabits}
\acro{GB}{Gigabytes}
\acro{GbE}{Gigabit Ethernet}
\acro{GigE}{Gigabit Ethernet}
\acro{GMT}{Giant Magellan Telescope}
\acro{GMVA}{Global Millimeter VLBI Array}
\acro{GOUS}{Group Observing Unit Set}
\acro{GPS}{Global Positioning System}
\acro{GPU}{Graphics Processing Unit}
\acro{GSps}{Giga-samples per second}
\acro{GUI}{Graphical User Interface}
\acro{HDD}{Hard Disk Drive}
\acro{HEMT}{High-electron-mobility Transistor}
\acro{HILSE}{Hardware in the Loop Simulation Environment}
\acro{HPBW}{Half Power Beam Width}
\acro{HVAC}{Heating, Ventilation, and Air Conditioning}
\acro{HW}{Hardware}
\acro{I/F}{Interface}
\acro{IAS}{Integrated Alarm System}
\acro{ICD}{Interface Control Document}
%\acro{ICM}{Imaging Correlation Mode} No such mode
\acro{ICT}{Integrated Computing Team}
\acro{ICV}{Imaging Correlation (with or without) VLBI Beamforming}
\acro{ICZoom}{Imaging Correlation (with or without) Zoom}
\acro{IDT}{Integrated Development Team}
\acro{IDTF}{Integrated Development Test Facility}
\acro{IET}{Integrated Engineering Team}
%\acro{IF}{Interferometric}
\acro{IF}{Intermediate Frequency}
\acro{IF-PL}{Interferometric Pipeline}
\acro{IFP}{Intermediate Frequency Processor}
\acro{IFSC}{IF Switch/Conditioner}
\acro{IFSCM}{IF Switch/Conditioner Module}
\acro{IMG}{Infrastructure Maintenance Group}
\acro{IMOP}{Integrated Maintenance and Operation Plan}
\acro{IRM}{Integration Release Manager}
\acro{iSCSI}{Internet Small Computer System Interface}
\acro{ISOpT}{Integrated Science Operations Team}
\acro{IST}{Integrated Science Team}
\acro{IT}{Information Technology}
\acro{IWS}{Initial WSU System}
\acro{IxT}{Integrated x Team}
\acro{JAO}{Joint ALMA Observatory}
\acro{JBOD}{Just a Bunch Of Disks}
\acro{JDT}{JAO Development Team}
\acro{KPI}{Key Performance Indicator}
\acro{LAB}{Laboratoire d'Astrophysique de Bordeaux}
\acro{LLC}{Line Length Corrector}
\acro{LNA}{Low Noise Amplifier}
\acro{LO}{Local Oscillator}
\acro{LO1}{First Local Oscillator}
\acro{LO2}{Secondary Local Oscillator}
\acro{LORR}{Local Oscillator Reference Receiver}
\acro{LORTM}{Local Oscillator Reference Test Module}
\acro{LPR}{Local Oscillator Photonic Receiver}
\acro{LPRA}{LO Photonic Reference Article}
\acro{LRU}{Line Replaceable Unit}
\acro{LS}{Laser Synthesizer}
\acro{LSB}{Lower Sideband}
\acro{LST}{Local Sidereal Time}
\acro{LVDS}{Low Voltage Differential Signaling}
\acro{MACI}{Management and Control Interface}
\acro{MC}{Monitor and Control}
\acro{MCF}{Monitor Control and Formatting}
\acro{MD}{Mixer-Detuning}
\acro{MDR}{Manual Data Reduction}
\acro{MDS}{Monitoring Data System}
\acro{MELCO}{Mitsubishi Electric Corporation}
\acro{MFS}{Master Frequency Standard}
%\acro{MFS}{Multi-frequency Synthesis}
\acro{ML}{Master Laser}
\acro{MOUS}{Member Observing Unit Set}
\acro{MPI}{Message Passing Interface}
\acro{MPLS}{Multiprotocol Label Switching}
\acro{MQTT}{Message Queueing Telemetry Transport}
\acro{MRR}{Manufacturing Readiness Review}
\acro{MS}{Measurement Set}
\acro{MSv2}{MS version 2}
\acro{MSv4}{MS version 4}
\acro{MTTF}{Mean Time To Failure}
\acro{MW}{MegaWatt}
\acro{MWA}{Murchison Widefield Array}
\acro{MWS}{Minimum WSU System}
\acro{N/A}{Not Applicable}
\acro{NA}{North American}
\acro{NAOJ}{National Astronomical Observatory of Japan}
\acro{NCR}{Non-Conformance Report}
\acro{NED}{NASA/IPAC Extragalactic Database}
\acro{NGAS}{Next Generation Archive System}
%\acro{ngCASA}{New Generation CASA}
%\acro{ngDPS}{New Generation Data Processing System}
\acro{ngFEMC}{Next Generation Front-End Monitor and Control System}
\acro{ngVLA}{Next Generation Very Large Array}
\acro{NRAO}{National Radio Astronomy Observatory}
\acro{NRC}{National Research Council Canada}
\acro{NTP}{Network Time Protocol}
\acro{NVME}{Non-Volatile Memory Express}
\acro{OCRO}{Operations Support Facility Correlator Room}
\acro{OIF}{Optical Internetworking Forum}
\acro{OMC}{Operator Monitoring and Control}
\acro{OOK}{On-Off Keying}
\acro{OPC/UA}{OPC Unified Architecture}
\acro{OPEX}{Operating Expenses or Expenditure}
\acro{OSA}{Optical Spectrum Analyzer}
\acro{OSF}{Operations Support Facility}
\acro{OSNR}{Optical Signal-to-Noise Ratio}
\acro{OSPFB}{Oversampling Polyphase Filter Bank}
\acro{OT}{Observing Tool}
\acro{OTDR}{Optical Time Domain Reflectometer}
\acro{OTF}{on-the-fly}
\acro{OUS}{Observing Unit Set}
\acro{P2G}{Phase 2 Generation}
\acro{PA}{Product Assurance}
\acro{PAI}{Preliminary Acceptance In-house}
\acro{PAS}{Provisional Acceptance on Site}
\acro{PB}{Petabytes}
%\acro{PB}{Primary Beam}
\acro{PCB}{Printed Circuit Board}
\acro{PDF}{Portable Document Format}
\acro{PDR}{Preliminary Design Review}
\acro{PDU}{Power Distribution Unit}
\acro{Ph1M}{Phase 1 Manager}
\acro{PHT}{Proposal Handling Team}
\acro{PI}{Principal Investigator}
\acro{PL}{Pipeline}
\acro{PLL}{Phase Lock Loop}
\acro{PLO}{Photonic Local Oscillator}
\acro{PLWG}{Pipeline Working Group}
\acro{PM}{Total Power array Mitsubishi antenna}
\acro{PPS}{Pulse Per Second}
\acro{PR}{Photonic Reference}
\acro{ProMo}{Product Mover}
\acro{ProTrack}{Project Tracker}
\acro{PRTSIR}{Problem Report and Tracking System - Investigation Report}
\acro{PRTSPR}{Problem Report and Tracking System - Problem Report}
\acro{PSA}{Power Supply Analog Rack LRU}
\acro{PSAD}{Power Supply Analog Distribution}
\acro{PSO}{Per Source Observing}
\acro{PTN}{Product Tree Number}
\acro{PTP}{Precision Time Protocol}
\acro{PU}{Processing Unit}
\acro{pwv}{precipitable water vapor; also styled as PWV}
\acro{QA}{Quality Assurance}
\acro{QAn}{Quality Assurance level n}
\acro{QB}{Queue Building}
\acro{QSFP}{Quad Small Formfactor Pluggable}
\acro{RADPS-ALMA}{Radio Astronomy Data Processing System - ALMA}
\acro{RADPS}{Radio Astronomy Data Processing System}
\acro{RAE}{Retrofitted Array Element}
\acrodef{ram-memory}[RAM]{Random Access Memory}
\acrodef{ram-reliability}[RAM]{Reliability, Availability, Maintainability}
\acro{RDMA}{Remote Direct Memory Acccess}
\acro{RF}{Radio Frequency}
\acro{RFI}{Radio Frequency Interference}
\acro{RH}{Request Handler}
\acro{RHEL}{Red Hat Enterprise Linux}
\acro{RID}{Review Item Discrepancy}
\acroplural{RID}[RIDs]{Review Item Discrepancies}
\acro{RISC}{Reduced Instruction Set Computer}
\acro{RMS}{Root Mean Square}
\acro{SACM}{Science Archive Content Manager}
\acro{SAS}{SubArray Switch (time and frequency distribution context)}
\acro{SB}{Scheduling Block}
\acro{SC}{Spectral Checkout}
\acro{SCCB}{Software Change Control Board}
\acro{SCO}{Santiago Central Office}
\acro{SD-PL}{Single-Dish Pipeline}
\acro{SD}{Single-Dish}
\acro{SDD}{System Design Description}
\acro{SE}{Systems Engineering}
\acro{SEG}{Systems Engineering Group}
\acro{SEMP}{Systems Engineering Management Plan}
\acro{SEU}{Single Event Upset}
\acro{SG}{Science Goal}
\acro{SGOUS}{Science Goal Observing Unit Set}
\acro{SIS}{Superconductor Insulator Superconductor}
\acro{SL}{Slave Laser}
\acro{SMA}{SubMiniature version A}
\acro{SNE}{Series Network Engine}
\acro{SnooPI}{Snooping Project Interface}
\acro{SNR}{Signal-to-Noise Ratio}
\acro{SONET}{Synchronous Optical Networking}
\acro{SOW}{Statement of Work}
\acro{SPC}{Signal Path Connectivity}
\acro{SPT}{South Pole Telescope}
\acro{spw}{Spectral Window}
\acro{SQLD}{Square Law Detector}
\acro{SRR}{System Requirements Review}
\acro{SSB}{Single Sideband}
\acro{SSO}{Solar System Object}
\acro{SSR}{Science Software Requirements}
\acro{SSS}{Subsystem Scientist}
\acro{STE}{Standard Test Environment}
\acro{SW}{Software}
\acro{TB}[TB]{Technical Building}
\acro{TBC}{To Be Confirmed (at a later date)}
\acro{TBD}{To Be Defined/Determined (at a later date)}
\acro{TBY}[TB]{Terabytes}
\acro{tCLOTS}{Temporary Central Local Oscillator Test Stand}
\acro{TCP/IP}{Transmission Control Protocol / Internet Protocol}
\acro{TE}{Timing Event}
\acro{TelCal}{Telescope Calibration}
\acro{TF}{Technical Facility}
\acro{TIA}{Telecommunications Industry Association}
\acro{TMC}{Telescope Monitor \& Control}
\acro{TMCDB}{Telescope Monitoring \& Configuration Data Base}
\acro{TMT}{Thirty Meter Telescope}
\acro{TP}{Total Power}
\acro{TPA}{Total Power Array}
\acro{TPGS}{Total Power GPU Spectrometer}
\acro{TSI}{Technical Solution Investigation}
\acro{Tsys}{System Temperature}
\acro{TWO}{Transition to WSU Operations}
\acro{UI}{User Interface}
\acro{UID}{Unique Identifier}
\acro{UPS}{Uninterruptible Power Supply}
\acro{URR}{Upgrade Readiness Review}
\acro{USB}{Upper Sideband}
\acro{UTC}{Coordinated Universal Time}
\acro{VAT}{Validation and Testing}
\acro{VB}{VLBI Beamforming}
\acro{VESDA}{Very Early Smoke Detection Apparatus}
\acro{VLAN}{Virtual Local Area Network}
\acro{VLASS}{Very Large Array Sky Survey}
\acro{VLBI}{Very Long Baseline Interferometry}
\acro{VSWR}{Voltage Standing Wave Ratio}
\acro{WCA}{Warm Cartridge Assembly}
\acro{WDM}{Wavelength Division Multiplexing}
\acro{WG}{Working Group}
\acro{WIFP}{Wideband Intermediate Frequency Processor}
\acro{WSU}{Wideband Sensitivity Upgrade}
\acro{WTS}{WSU Test System}
\acro{WVR}{Water Vapor Radiometer}
\acro{XML}{eXtensible Markup Language}
\acro{XTSS}{eXTended State System}
\acro{YTO}{YIG Tuned Oscillator}
\acro{ZF}{Zoom Factor}
%%%%%%%%%%%%%%%
\end{acronym}
%%%%%%%%%%%%%%%
%%%%%%%%%%%%%%%%%%%%%%%%%%%%%%%%%%